\subsubsection{Rotating calipers (polygon diameter)}
\begin{lstlisting}[style=mystyle, language=C++]
// TC: O(N) (assumes convex polygon sorted counter-clockwise)
// usa: encontrar el par de vertices mas distantes (diametro) en un poligono convexo
#include <bits/stdc++.h>
using namespace std;

long long rotatingCalipers(const vector<Point>& p){
	int n = (int)p.size();
	if(n == 1) return 0;
	if(n == 2){
		long long dx = p[0].x - p[1].x, dy = p[0].y - p[1].y;
		return dx*dx + dy*dy;
	}
	long long best = 0;
	int k = 1;
	for(int i=0;i<n;i++){
		int ni = (i + 1) % n;
		while(true){
			int nk = (k + 1) % n;
			long long cur = abs((long long)(p[ni] - p[i]) * (p[nk] - p[k]));
			long long nxt = abs((long long)(p[ni] - p[i]) * (p[(nk+1)%n] - p[(k+1)%n]));
			if(nxt > cur) k = nk;
			else break;
		}
		long long d1 = 1LL * (p[i].x - p[k].x) * (p[i].x - p[k].x) + 1LL * (p[i].y - p[k].y) * (p[i].y - p[k].y);
		long long d2 = 1LL * (p[ni].x - p[k].x) * (p[ni].x - p[k].x) + 1LL * (p[ni].y - p[k].y) * (p[ni].y - p[k].y);
		best = max(best, max(d1, d2));
	}
	return best;
}

int main(){
	vector<Point> tri = {{0,0}, {10,0}, {5,8}};  // triangle CCW
	cout << "diameter^2 = " << rotatingCalipers(tri) << "\n";
	return 0;
}
\end{lstlisting}
